% Compile with: pdflatex -shell-escape demo.tex
% (minted requires -shell-escape; pygments-rzk must be installed)
\documentclass[aspectratio=169]{beamer}

\usetheme{moloch}  % or: \usetheme{metropolis}

\usepackage{minted}
\usemintedstyle{xcode}

% Suppress the red error boxes Pygments draws around tokens the rzk lexer
% does not recognise (the lexer has no catch-all rule yet).
\AtBeginDocument{%
  \expandafter\def\csname PYG@tok@err\endcsname{}%
}

\title{\textsc{Rzk} highlighting in LaTeX with \texttt{minted}}
\author{}
\date{}

\begin{document}

\begin{frame}[fragile]
  \frametitle{\textsc{Rzk} highlighting in LaTeX with \texttt{minted}}

A basic example:

\scriptsize
\begin{minted}[linenos,frame=leftline,texcomments]{rzk}
#lang rzk-1

#section path-algebra

#variable A : U
#variables x y z : A

-- $\text{Paths are symmetric, so we can reverse $x =_A y$ to $y =_A x$.}$
#define rev uses (A x y)
  (p : x = y)       -- $\text{A path from $x$ to $y$ in $A$.}$
  : y = x           -- $\text{The reversal will be defined by path induction on $p$.}$
  := idJ(A, x, \y' p' -> y' = x, refl, y, p)

-- $\text{Path composition (concatenation) by induction on the second path.}$
#define concat
  (p : x = y)       -- $\text{A path from $x$ to $y$ in $A$.}$
  (q : y = z)       -- $\text{A path from $y$ to $z$ in $A$.}$
  : (x = z)         -- $\text{The concatenated path $x = y = z$.}$
  := idJ(A, y, \z' q' -> (x = z'), p, z, q)

#end path-algebra
\end{minted}

\end{frame}

\end{document}
